\section{Related Work and Positioning}
\label{sec:related}

% RW-P1 — Multi-hop / iterative RAG
Prior multi-hop and agentic RAG systems improve evidence acquisition by
iterating retrieval, reformulating queries, or traversing explicit structures.
Their central strength is adaptive access to later-hop evidence. However, the
information need to be satisfied and the physical actions used to satisfy it
are often coupled in a method-specific control flow, making equivalent evidence
goals difficult to compare and optimize across alternative executions. We do
not claim these systems cannot plan; rather, the evidence condition to be met
and the retrieval procedure that meets it are not separated in a way an
optimizer can exploit.

% RW-P2 — Planning and executable-program RAG
Recent work has made RAG increasingly explicit: PlanRAG organizes atomic queries
into logical query trees, while PyRAG exposes retrieval and reasoning as
executable programs. These advances improve inspectability and planning.
SlotRAG addresses a different systems question: how to represent the evidence
conditions that must hold independently of the physical retrieval strategy, and
how to optimize multiple physical realizations of the same evidence program
under runtime uncertainty. The contribution is not ``also executable''; it is
the separation of a typed, optimizer-visible evidence program from the physical
plan that realizes it, and the requirement-aware allocation over that
separation.

% RW-P3 — Learned adaptive retrieval/control
A complementary line learns policies for deciding when and how to retrieve from
an evolving state (DynaKRAG; active RAG; stopping policies). SlotRAG
deliberately separates prediction from optimization: learned or derived
estimators predict properties such as evidence yield, selectivity, failure
probability, and resource cost, while a deterministic allocator uses those
estimates to assign execution order, retrieval implementation, and budget to
each requirement. The separation makes the allocation auditable and the
estimator swappable; general plan re-enumeration and dominance pruning across
arbitrary graph topologies are future work under a branching execution
model.

% RW-P4 — Declarative semantic data systems
Declarative AI data systems have already demonstrated the value of semantic
operators, physical alternatives, cost estimation, and adaptive execution
(LOTUS, Palimpzest/Abacus, Sema). Their workloads are general semantic data
transformations. SlotRAG specializes these systems principles to dependent
evidence acquisition, where operator utility is defined by progress toward
unresolved evidence requirements and where provenance and evidence survival
into the answer context are first-class execution outcomes. The objective
(Equation~\ref{eq:objective}) is deliberately not a generic cost optimizer:
utility is credited only for marginal progress on requirements still
unsatisfied, which distinguishes it from cost or relevance optimization over
general transformation pipelines.

% RW-P5 — Positioning synthesis
Table~\ref{tab:positioning} positions SlotRAG along four dimensions:
declarative evidence requirements, a typed evidence algebra, logical/physical
separation, and requirement-aware physical planning. The contribution is not any
isolated primitive---planning, executable programs, learned estimates, semantic
operators, and adaptive execution all exist---but an evidence-centric execution
model that jointly provides these under explicit budgets, with evidence state
as the interface between the declarative program and physical optimization.

\begin{table}[t]
\centering
\caption{Positioning of SlotRAG against closest work along four dimensions.
\checkmark\ = first-class/central; $\circ$ = partial/peripheral; --- = absent.
Positioning is by design/method only; none of the listed systems (PlanRAG,
PyRAG, LOTUS, Sema) was re-run as an experimental baseline in this paper---they
are cited as related work, not as evaluated competitors.}
\label{tab:positioning}
\small
\begin{tabular}{lccccc}
\toprule
 & SlotRAG & PlanRAG & PyRAG & LOTUS & Sema \\
\midrule
Declarative evidence requirements & \checkmark & $\circ$ & $\circ$ & --- & $\circ$ \\
Typed evidence algebra w/ state transitions & \checkmark & --- & --- & $\circ$ & --- \\
Logical/physical separation & \checkmark & $\circ$ & --- & \checkmark & \checkmark \\
Requirement-aware budgeted planning & \checkmark & $\circ$ & --- & --- & $\circ$ \\
Provenance as first-class outcome & \checkmark & --- & --- & --- & --- \\
\bottomrule
\end{tabular}
\end{table}
